\chapter{Literature Review}
\label{chap:lr}
\chaptermark{Second Chapter Heading}

\section{Language Server Protocol}
\label{lr:lsp}

The Language Server Protocol developed by Microsoft\cite{lspspecification},
has profoundly influenced the development of both domain-specific languages\footnote{Commonly abbreviated as DSLs}
and integrated development environments by decoupling them via a standardized set of requests.
Consequently, DSL developers no longer need to implement and maintain language-specific logic for each supported IDE;
conversely, IDE integration requires significantly less effort by leveraging these standardized API calls.

After studying the impact of LSP in \cite{lspimpact}, Hendrik B{\"u}ender concludes that
the introduction of the protocol not only substantially reduced development effort,
but also defined \enquote{minimum feature set that should be supported by any DSL editor}
and even enabled the complete omission of local IDE installation in certain scenarios.

\section{Rzk}
\label{lr:rzk}

Rzk is the interactive theorem prover developed \enquote{with the idea of bringing Riehl and Shulman's
paper\cite{riehl2023typetheorysyntheticinftycategories} to life by implementing a proof assistant
based on their type thory with shapes}, as stated in its official GitHub repository\cite{rzkcompilercodebase}.

It has already been utilized in several mathematical formalization projects:
\begin{itemize}
	\item Formalizing the $\infty$-Categorical Yoneda Lemma \cite{yonedalemma}
	\item Simplicial HoTT and synthetic $\infty$-categories \cite{shott}
\end{itemize}

\section{Typechecker as a Build System}
\label{lr:buildsys}

In \cite{lacarte}, Mokhov et al. propose a systematic framework for developing build systems.
This framework conceptualizes a build system as two architectural components:
a \enquote{Scheduler}, handles task scheduling in dependency order,
and a \enquote{Rebuilder}, which determines whether a task requires re-execution.

The Rzk language server compiles a project into a collection of diagnostics (errors and warnings),
storing an intermediate representation that contains declarations and diagnostics for each module.
By treating project files as \enquote{inputs}, module caches as \enquote{outputs},
and the typechecking process itself as the primary task,
the aforementioned framework can be leveraged to develop an incremental build system.
This approach eliminates repetitive checking of unmodified code,
thereby lowering latency under specific conditions.

Other notable frameworks exist within the domain of incremental computation systems:
\begin{itemize}
	\item In \cite{adapton}, Hammer et al. build upon the concept of \enquote{demanded computations graph}
	to present theoretical \enquote{core calculus that applies a demand-driven semantics to incremental computation},
	implementing this calculus as an OCaml library named Adapton.
	\item The rust-analyzer team developed Salsa\cite{salsa}, a Rust framework
	that utilizes a similar incremental, demand-driven computation.
\end{itemize}

Despite both framework introducing powerful paradigms for implementing incremental computation,
they are fundamentally based on the demand-driven paradigm\cite{reportonsalsaadapton}.
This contrasts with Rzk, which enforces a strict project structure and requires full compilation.
Furthermore, it is noteworthy that both Rzk and framework presented in \cite{lacarte} are implemented in Haskell,
whereas Adapton is an OCaml library and Salsa is written in Rust.

In conclusion, the build system framework proposed by Mokhov et al.
is more aligned with the objectives of this thesis,
though specific architectural details from Adapton and Salsa may still influence the final implementation.

\section{Document processing in various ITPs}
\label{lr:itps}

\subsection{Roqc}

Roqc Prover (formerly known as Coq) is an interactive theorem prover
based on the Calculus of Inductive Constructions.

Prior to the introduction of the LSP, Roqc utilized a proprietary implementation of
asynchronous document processing, as described in \cite{barras2015asynchronousprocessingcoqdocuments}.
That work demonstrates significant latency reduction both during
the batch compilation of entire project and during on-edit evaluation within CoqIDE.
However, that approach focuses on parallel and asyncronous processing
of independent documents rather than on incremental computation of unified project.

As noted in its package documentation\cite{coq-lsp}, the Rocq LSP server (coq-lsp)
compiles code incrementally and continuously validates documents as user edits arrive.

Rocq compiles its '.v' source code files into '.vo' object files.
Because the feedback loop in IDE is based on importing these pre-compiled object files,
the process inherently exhibits incremental characteristics.

% \subsection{Lean4}
%
% Lean is a functional programming language, theorem prover, and a metaprogramming environment.
% Its language server re-evaluates code upon each keystroke,
% cacheing both individual declarations within a file and the compiled states of individual modules.
%
% \subsection{Agda}
%
% Agda is a dependently typed programming language and an extension of Martin-Löf’s type theory.
% It evaluates only when a file is saved or when an evaluation is manually requested.
% While Agda caches compilation results at the file level, any subsequent modification
% to a file forces a complete re-evaluation of that entire file.
%
% \subsection{Rzk}
%
% \subsection{Improved Rzk}

\subsection{Summary}
To conclusde this section, the capabilities and behaviors
of these varius environments are summarized in Table\ref{tab:feats}.
\begin{table}[b]
	\centering
	\begin{center}
		\begin{tabular}{|c|c|c|c|c|c|}
			\hline
			Language server & Reevaluation trigger & Project-level cache & Cache structure \\ \hline
			Rocq & Edit & Compiled modules & Module imports \\ \hline
			% Lean4 & Edit & Compiled modules & DAG\tablefootnote{Directed acyclic graph} \\ \hline
			% Agda & Request/save & Compiled modules & Module imports \\ \hline
			% Rzk & Save & Typed module caches & Flat list \\ \hline
			% Improved Rzk & Edit & Typed module caches & TBD \\ \hline
		\end{tabular}
	\end{center}
	\caption{Caching and evaluation triggers in ITPs}
	\label{tab:feats}
\end{table}
